\section*{Appendices}
\addcontentsline{toc}{section}{\emph{Appendices}}

\section{Contributions}

\begin{table}[H]
\centering
\caption{Individual contributions of each authors}
\label{tab:contributions}
\begin{tabular}{@{}lll@{}}
\toprule
\textbf{Name} & \textbf{Implication} & \textbf{Main contributions} \\ \midrule
Simone Binotto & 95 & Resistojets, Electrothermal thrusters\\
Marco Azevedo & 100 & Arcjets, SW Sensitivity analysis \\
Philippe Quichaud & 100 & Pulsed plasma thrusters, SW Scaling laws \\
Corentin Ouisse & 105 & Hall effect thrusters, SW User interface  \\
Andr\'{e} Matos & 100 & Ion thrusters, SW Scaling and sweeping \\
Th\'{e}o Demesy & 95 & Introduction / Conclusion, Review \\ \bottomrule
\end{tabular}
\end{table}

\section{Software}

The team developped, in addition to this report, a software written in Python that allows the user to dimension an Ion thruster for a given mission. The user can freely specify mission parameters, as for example desired delta-v and payload mass in a text file, but also more specific parameters relative to precise dimensionning of the motor (minimum thrust, component efficiencies, and so on...).

This text file will then be processed by the program which will return an optimal motor design (optimized for mass) described in an output file. It will also provide a table all possible motor configurations encountered in the search process. Finally, it also conducts a sensibility analysis on a specified design objective (eg mass) relative to all mission parameters so that the user can evaluate the effect of design choices.

For further information relative to the software, please refer to the provided README file that explains in detail usage and features.